\begin{proof}[Proof of \cref{obj:thm:stable-grid-upper}]
Fix an experiment in \(\mathcal M_T^{(2)}(t_0,\zeta)\). All expectations below are under its observed trajectory law, denoted by \(\mathbb E_M\). Use the moments of \cref{obj:def:observable-moments} and the joint-state quantities of \cref{obj:lem:observable-intervention-moments}.

\begin{enumerate}
\item Use zero-based epochs throughout this proof: epoch \(t\) here corresponds to epoch \(t+1\) in \cref{obj:def:observable-moments}. Accordingly, define the shifted scores by
\[
Z_t(k):=Y_{t+1}\prod_{j=t-k+1}^{t+1}\rho_j,
\qquad 0\le k\le6,\quad k\le t<T,
\]
where the reward and likelihood ratios on the right retain their one-based indexing. Thus \(Z_t(k)\) is the score at one-based epoch \(t+1\), agreeing with the zero-based convention of \cref{obj:lem:observable-intervention-moments}.

Define the retained epoch set, its cardinality, and the maximum empirical moment error by
\[
\mathcal I_T:=\{6,\ldots,T-1\},
\qquad
n_T:=|\mathcal I_T|=T-6,
\qquad
\epsilon_T(w):=
\max_{0\le k\le6}|\widehat m_k(w)-m_k|,
\]
where \(w\in(\mathcal X\times\mathcal A\times\mathbb R)^T\) is an observed trajectory. Changing variables in the one-based retained-window sum gives
\[
\widehat m_k
=
\frac1{T-6}\sum_{t=6}^{T-1}Z_t(k)
=
\frac1{n_T}\sum_{t\in\mathcal I_T}Z_t(k),
\qquad 0\le k\le6.
\]
The parameter conditions give
\[
0<\alpha=\exp(-1/t_0)<1,
\qquad
L=\exp(\zeta)\ge1,
\qquad
B_\alpha\ge0,
\qquad
V_{\alpha,L}\ge0,
\]
and the horizon condition gives
\[
T>0,\qquad n_T>0,\qquad T\le2n_T,\qquad T\ge1.
\]
The maximum defining \(\epsilon_T(w)\) is attained among seven indices. Its square is therefore one of the nonnegative summands on the right-hand side of
\[
\epsilon_T(w)^2
\le
\sum_{k=0}^{6}(\widehat m_k(w)-m_k)^2.
\]
The observed law is a probability law, and the empirical moments are square-integrable. Thus the sum on the right is integrable; measurability of the finite maximum and the displayed domination also give integrability of \(\epsilon_T^2\).
% lean: CausalSmith.Stat.PomdpBinaryhiddenRate.seven_error_max_sq_le_sum
% lean: CausalSmith.Stat.PomdpBinaryhiddenRate.stableGrid_upper

\item For \(0\le k\le6\) and \(t,t'\in\mathcal I_T\), define
\[
c_{k;t,t'}
:=
\mathbb E_M\!\left[
  (Z_t(k)-m_k)(Z_{t'}(k)-m_k)
\right].
\]
Every retained epoch satisfies \(k\le t\). Consequently, \cref{obj:lem:observable-intervention-moments} supplies
\[
\mathbb E_M[Z_t(k)]=m_k=d_bP_e^kr_e,
\qquad
\mathbb E_M[Z_t(k)^2]\le L^{k+1}.
\]
In particular,
\[
0\le c_{k;t,t}
=
\mathbb E_M[(Z_t(k)-m_k)^2]
=
\mathbb E_M[Z_t(k)^2]-m_k^2
\le L^{k+1}.
\]
For either choice of sign, nonnegativity of the centered square gives
\[
\begin{aligned}
0
&\le
\mathbb E_M\!\left[
 \bigl((Z_t(k)-m_k)\pm(Z_{t'}(k)-m_k)\bigr)^2
\right]\\
&=
c_{k;t,t}+c_{k;t',t'}\pm2c_{k;t,t'}.
\end{aligned}
\]
Using both signs yields
\[
|c_{k;t,t'}|
\le
\frac{c_{k;t,t}+c_{k;t',t'}}2
\le L^{k+1}.
\]
For separated epochs, the covariance conclusion of
\cref{obj:lem:observable-intervention-moments} additionally gives
\[
|c_{k;t,t'}|
\le\alpha^{\,t'-t-k-1}
\qquad\text{when }t'-t\ge k+1.
\]
These identities and bounds apply under the observed law because the scores are functions of the observed trajectory.
% lean: CausalSmith.Stat.PomdpBinaryhiddenRate.observable_score_variance_le
% lean: CausalSmith.Stat.PomdpBinaryhiddenRate.observable_score_covariance_le
% lean: CausalSmith.Stat.PomdpBinaryhiddenRate.observable_score_covariance_le_separated

\item Fix \(0\le k\le6\) and \(t\in\mathcal I_T\). Partition its future retained epochs according to whether \(1\le t'-t\le k\) or \(t'-t\ge k+1\). There are at most \(k\) epochs in the first part. In the second part, the exponents \(t'-t-k-1\) are distinct integers in \(\{0,\ldots,T-1\}\). Hence
\[
\begin{aligned}
\sum_{\substack{t'\in\mathcal I_T\\t<t'}}
 |c_{k;t,t'}|
&=
\sum_{\substack{t'\in\mathcal I_T\\1\le t'-t\le k}}
 |c_{k;t,t'}|
+
\sum_{\substack{t'\in\mathcal I_T\\t'-t\ge k+1}}
 |c_{k;t,t'}|\\
&\le
kL^{k+1}+\sum_{j=0}^{T-1}\alpha^j\\
&=
kL^{k+1}+\frac{1-\alpha^T}{1-\alpha}\\
&\le
kL^{k+1}+\frac1{1-\alpha}.
\end{aligned}
\]
For \(k=0\), the first sum is empty and equals zero. Symmetry of covariance allows the full absolute covariance sum to be split into its diagonal and twice its upper triangle:
\[
\begin{aligned}
\sum_{t\in\mathcal I_T}\sum_{t'\in\mathcal I_T}c_{k;t,t'}
&\le
\sum_{t\in\mathcal I_T}\sum_{t'\in\mathcal I_T}|c_{k;t,t'}|\\
&=
\sum_{t\in\mathcal I_T}|c_{k;t,t}|
+
2\sum_{t\in\mathcal I_T}
 \sum_{\substack{t'\in\mathcal I_T\\t<t'}}
 |c_{k;t,t'}|\\
&\le
n_TL^{k+1}
+
2n_T\left(kL^{k+1}+\frac1{1-\alpha}\right)\\
&=
n_T\left((1+2k)L^{k+1}+\frac2{1-\alpha}\right)\\
&\le
n_T\left(13L^7+\frac2{1-\alpha}\right)
=
n_TV_{\alpha,L}.
\end{aligned}
\]
The last inequality uses \(k\le6\) and \(L\ge1\).

Averaging the score means over the common window gives
\[
\mathbb E_M[\widehat m_k]
=
\frac1{n_T}\sum_{t\in\mathcal I_T}\mathbb E_M[Z_t(k)]
=
m_k.
\]
Expanding the centered square of this finite average therefore yields
\[
\begin{aligned}
\mathbb E_M[(\widehat m_k-m_k)^2]
&=
\frac1{n_T^2}
\sum_{t\in\mathcal I_T}\sum_{t'\in\mathcal I_T}c_{k;t,t'}\\
&\le
\frac{V_{\alpha,L}}{n_T}.
\end{aligned}
\]
% lean: CausalSmith.Stat.PomdpBinaryhiddenRate.binary_sum_abs_covariance_future_le
% lean: CausalSmith.Stat.PomdpBinaryhiddenRate.binary_sum_abs_symmetric_eq_diag_add_two_upper
% lean: CausalSmith.Stat.PomdpBinaryhiddenRate.empiricalMoment_mse_eq_covariance_sum
% lean: CausalSmith.Stat.PomdpBinaryhiddenRate.empiricalMoment_mse_le_varianceFactor

\item Integrating \(\epsilon_T(w)^2\le\sum_{k=0}^{6}(\widehat m_k(w)-m_k)^2\) and applying the seven moment bounds gives
\[
\mathbb E_M[\epsilon_T^2]
\le
\sum_{k=0}^{6}\mathbb E_M[(\widehat m_k-m_k)^2]
\le
\frac{7V_{\alpha,L}}{n_T}.
\]
% lean: CausalSmith.Stat.PomdpBinaryhiddenRate.empiricalMoment_max_mse_le_varianceFactor

\item We next obtain a value bound for the selected grid realization. To distinguish the grid resolution from the fixed experiment, write
\[
M_{\mathrm{grid}}
:=
\left\lceil(22+36/\alpha)T\right\rceil,
\qquad
U:=\frac{\mathbf1\mathbf1^{\mathsf T}}4,
\qquad
\lambda:=\frac{\alpha}{\alpha+6/M_{\mathrm{grid}}}.
\]
Here \(\mathbf1\) is the four-dimensional vector of ones. Thus
\[
M_{\mathrm{grid}}>0,
\qquad
0<\lambda<1,
\qquad
P(R)=\lambda R+(1-\lambda)U.
\]
Use a fixed ordering of the four joint states as the four grid coordinates; this relabeling preserves scalar matrix moments and total variation.

The candidate set of \cref{obj:def:stable-grid-estimator} is finite and nonempty. Indeed,
\[
\bigl((1,0,0,0),\ \mathbf1(1,0,0,0),\ (0,0,0,0)\bigr)
\in\mathcal C_{M_{\mathrm{grid}}},
\]
because all its coordinates lie on the required grids and
\[
\delta\bigl(\mathbf1(1,0,0,0)\bigr)=0
\le\alpha+\frac6{M_{\mathrm{grid}}}.
\]
The lexicographic minimizer therefore belongs to this candidate set.

For every candidate \((\nu,R,r)\), its stabilized matrix is row-stochastic and satisfies
\[
P(R)_{ij}\ge\frac{1-\lambda}{4}.
\]
The common uniform component cancels between rows, so
\[
\begin{aligned}
\delta(P(R))
&=
\lambda\delta(R)\\
&\le
\lambda\left(\alpha+\frac6{M_{\mathrm{grid}}}\right)
=
\alpha.
\end{aligned}
\]
The same uniform minorization gives, for probability rows \(\mu,\mu'\),
\[
\|\mu P(R)-\mu'P(R)\|_{\mathrm{TV}}
\le
\lambda\|\mu-\mu'\|_{\mathrm{TV}}.
\]
Since \(\lambda<1\), contraction on the finite probability simplex supplies a stationary probability row:
\[
\pi(P(R))P(R)=\pi(P(R)),
\qquad
\pi(P(R))_i\ge0,
\qquad
\sum_{i=1}^{4}\pi(P(R))_i=1.
\]

Write the selected quantities as in \cref{obj:def:stable-grid-estimator}:
\[
(\nu_*,R_*,r_*)\in\mathcal C_{M_{\mathrm{grid}}},
\qquad
P_*:=P(R_*),
\qquad
\widehat\theta_T=\pi(P_*)r_*.
\]
The experiment's class conditions give a row-stochastic \(P_e\), probability rows \(d_b,d_e\), and a unit-range reward vector:
\[
P_e(i,j)\ge0,\quad \sum_jP_e(i,j)=1,
\qquad
d_b(i),d_e(i)\ge0,\quad \sum_id_b(i)=\sum_id_e(i)=1,
\]
\[
d_eP_e=d_e,
\qquad
0\le r_e(i)\le1.
\]
Applying \cref{obj:ass:target-contraction} to point masses gives
\[
\delta(P_e)
=
\max_{i,j}\frac12\sum_s|P_e(i,s)-P_e(j,s)|
\le\alpha.
\]
Together with \(\theta=d_er_e\) from \cref{obj:def:target-value}, these facts discharge the hypotheses of \cref{obj:lem:pair-polynomial-modulus}. Its value bound is
\[
\begin{aligned}
|\widehat\theta_T-\theta|
&\le
B_\alpha
\max_{0\le k\le6}
|\nu_*P_*^kr_*-d_bP_e^kr_e|\\
&=
B_\alpha
\max_{0\le k\le6}
|\nu_*P_*^kr_*-m_k|.
\end{aligned}
\]
% lean: CausalSmith.Stat.PomdpBinaryhiddenRate.selectGridCode_mem
% lean: CausalSmith.Stat.PomdpBinaryhiddenRate.stabilizedGridMatrix_stationary
% lean: CausalSmith.Stat.PomdpBinaryhiddenRate.stabilizedGridMatrix_dobrushin_le
% lean: CausalSmith.Stat.PomdpBinaryhiddenRate.grid_value_modulus

\item Construct an approximating candidate by rounding. For every probability row \(\mu\) on the four coordinates, define
\[
\bigl(\mathcal R_{\mathrm{grid}}(\mu)\bigr)_i
:=
\begin{cases}
\lfloor M_{\mathrm{grid}}\mu_i\rfloor/M_{\mathrm{grid}},
 &i=1,2,3,\\[2mm]
1-\displaystyle\sum_{j=1}^{3}
 \lfloor M_{\mathrm{grid}}\mu_j\rfloor/M_{\mathrm{grid}},
 &i=4.
\end{cases}
\]
The first three coordinates are rounded down, and their sum is at most one. The fourth coordinate is therefore nonnegative, is a grid multiple, and makes the total mass one. For \(i=1,2,3\),
\[
0\le
\mu_i-\bigl(\mathcal R_{\mathrm{grid}}(\mu)\bigr)_i
\le\frac1{M_{\mathrm{grid}}},
\]
while
\[
\mu_4-\bigl(\mathcal R_{\mathrm{grid}}(\mu)\bigr)_4
=
-\sum_{i=1}^{3}
 \left(\mu_i-\bigl(\mathcal R_{\mathrm{grid}}(\mu)\bigr)_i\right).
\]
Consequently,
\[
\begin{aligned}
\|\mu-\mathcal R_{\mathrm{grid}}(\mu)\|_{\mathrm{TV}}
&=
\frac12\sum_{i=1}^{4}
 \left|\mu_i-\bigl(\mathcal R_{\mathrm{grid}}(\mu)\bigr)_i\right|\\
&=
\sum_{i=1}^{3}
 \left(\mu_i-\bigl(\mathcal R_{\mathrm{grid}}(\mu)\bigr)_i\right)
\le\frac3{M_{\mathrm{grid}}}.
\end{aligned}
\]

Define the rounded initial row, transition rows, and rewards by
\[
\nu_{\mathrm{grid}}:=\mathcal R_{\mathrm{grid}}(d_b),
\qquad
R_{\mathrm{grid}}(i,\cdot)
:=
\mathcal R_{\mathrm{grid}}(P_e(i,\cdot)),
\qquad
r_{\mathrm{grid}}(i)
:=
\frac{\lfloor M_{\mathrm{grid}}r_e(i)\rfloor}{M_{\mathrm{grid}}}.
\]
They satisfy
\[
\|\nu_{\mathrm{grid}}-d_b\|_{\mathrm{TV}}
\le\frac3{M_{\mathrm{grid}}},
\qquad
\max_i\|R_{\mathrm{grid}}(i,\cdot)-P_e(i,\cdot)\|_{\mathrm{TV}}
\le\frac3{M_{\mathrm{grid}}},
\]
\[
0\le r_{\mathrm{grid}}(i)\le1,
\qquad
0\le r_e(i)-r_{\mathrm{grid}}(i)
\le\frac1{M_{\mathrm{grid}}}.
\]
For each pair of rows, the triangle inequality gives
\[
\begin{aligned}
\|R_{\mathrm{grid}}(i,\cdot)-R_{\mathrm{grid}}(j,\cdot)\|_{\mathrm{TV}}
&\le
\|R_{\mathrm{grid}}(i,\cdot)-P_e(i,\cdot)\|_{\mathrm{TV}}\\
&\quad+
\|P_e(i,\cdot)-P_e(j,\cdot)\|_{\mathrm{TV}}\\
&\quad+
\|P_e(j,\cdot)-R_{\mathrm{grid}}(j,\cdot)\|_{\mathrm{TV}}\\
&\le\alpha+\frac6{M_{\mathrm{grid}}}.
\end{aligned}
\]
Taking the maximum over the row pair proves
\[
\delta(R_{\mathrm{grid}})
\le\alpha+\frac6{M_{\mathrm{grid}}},
\qquad
(\nu_{\mathrm{grid}},R_{\mathrm{grid}},r_{\mathrm{grid}})
\in\mathcal C_{M_{\mathrm{grid}}}.
\]
% lean: CausalSmith.Stat.PomdpBinaryhiddenRate.reward_grid_rounding
% lean: CausalSmith.Stat.PomdpBinaryhiddenRate.simplex_grid_rounding
% lean: CausalSmith.Stat.PomdpBinaryhiddenRate.gridDobrushin_perturbation
% lean: CausalSmith.Stat.PomdpBinaryhiddenRate.grid_candidate_moment_approximation

\item Define the stabilized approximating matrix by
\[
P_{\mathrm{grid}}
:=
P(R_{\mathrm{grid}})
=
\lambda R_{\mathrm{grid}}+(1-\lambda)U.
\]
For every row, its distance from the uniform probability row is at most one. Also,
\[
1-\lambda
=
\frac{6/M_{\mathrm{grid}}}{\alpha+6/M_{\mathrm{grid}}}
\le\frac6{\alpha M_{\mathrm{grid}}}.
\]
It follows that
\[
\begin{aligned}
\|R_{\mathrm{grid}}(i,\cdot)-P_{\mathrm{grid}}(i,\cdot)\|_{\mathrm{TV}}
&=
(1-\lambda)
\|R_{\mathrm{grid}}(i,\cdot)-U(i,\cdot)\|_{\mathrm{TV}}\\
&\le\frac6{\alpha M_{\mathrm{grid}}},
\end{aligned}
\]
and hence
\[
\|P_e(i,\cdot)-P_{\mathrm{grid}}(i,\cdot)\|_{\mathrm{TV}}
\le
\frac{3+6/\alpha}{M_{\mathrm{grid}}}.
\]

For every integer \(k\ge0\), stochasticity and the reward bounds imply
\[
0\le(P_e^kr_e)(i)\le1.
\]
Integration of a \([0,1]\)-valued function against the difference of two probability rows is bounded by their total variation distance: centering that function at \(1/2\) gives this bound directly. We claim
\[
\|P_e^kr_e-P_{\mathrm{grid}}^kr_{\mathrm{grid}}\|_\infty
\le
\frac{k(3+6/\alpha)+1}{M_{\mathrm{grid}}},
\qquad k\ge0.
\]
For \(k=0\), this is the reward-rounding bound. At the induction step, insert
\(P_{\mathrm{grid}}P_e^kr_e\) between the two expressions. The row approximation bound controls the first difference, and averaging the previous error with a stochastic row controls the second:
\[
\begin{aligned}
&\|P_e^{k+1}r_e-P_{\mathrm{grid}}^{k+1}r_{\mathrm{grid}}\|_\infty\\
&\quad\le
\|(P_e-P_{\mathrm{grid}})P_e^kr_e\|_\infty
+
\|P_{\mathrm{grid}}
 (P_e^kr_e-P_{\mathrm{grid}}^kr_{\mathrm{grid}})\|_\infty\\
&\quad\le
\frac{3+6/\alpha}{M_{\mathrm{grid}}}
+
\|P_e^kr_e-P_{\mathrm{grid}}^kr_{\mathrm{grid}}\|_\infty\\
&\quad\le
\frac{(k+1)(3+6/\alpha)+1}{M_{\mathrm{grid}}}.
\end{aligned}
\]
The initial-distribution error can now be separated from the power-and-reward error:
\[
\begin{aligned}
|m_k-\nu_{\mathrm{grid}}P_{\mathrm{grid}}^kr_{\mathrm{grid}}|
&=
|d_bP_e^kr_e-\nu_{\mathrm{grid}}P_{\mathrm{grid}}^kr_{\mathrm{grid}}|\\
&\le
|(d_b-\nu_{\mathrm{grid}})P_e^kr_e|
+
|\nu_{\mathrm{grid}}
 (P_e^kr_e-P_{\mathrm{grid}}^kr_{\mathrm{grid}})|\\
&\le
\frac3{M_{\mathrm{grid}}}
+
\frac{k(3+6/\alpha)+1}{M_{\mathrm{grid}}}\\
&=
\frac{4+k(3+6/\alpha)}{M_{\mathrm{grid}}}.
\end{aligned}
\]
For \(0\le k\le6\), the coefficient and ceiling calibration give
\[
4+k(3+6/\alpha)\le22+36/\alpha,
\qquad
(22+36/\alpha)T\le M_{\mathrm{grid}},
\]
so
\[
|m_k-\nu_{\mathrm{grid}}P_{\mathrm{grid}}^kr_{\mathrm{grid}}|
\le
\frac{22+36/\alpha}{M_{\mathrm{grid}}}
\le\frac1T.
\]
% lean: CausalSmith.Stat.PomdpBinaryhiddenRate.stabilizedGridMatrix_row_error
% lean: Causalean.Mathlib.Probability.FiniteMarkovPerturbation.stochastic_power_reward_error
% lean: Causalean.Mathlib.Probability.FiniteMarkovPerturbation.stochastic_moment_error
% lean: CausalSmith.Stat.PomdpBinaryhiddenRate.grid_approximation

\item Fix an observed trajectory \(w\). The approximating candidate has empirical discrepancy bounded by
\[
\begin{aligned}
\max_{0\le k\le6}
|\widehat m_k(w)-\nu_{\mathrm{grid}}P_{\mathrm{grid}}^kr_{\mathrm{grid}}|
&\le
\max_{0\le k\le6}|\widehat m_k(w)-m_k|\\
&\quad+
\max_{0\le k\le6}
|m_k-\nu_{\mathrm{grid}}P_{\mathrm{grid}}^kr_{\mathrm{grid}}|\\
&\le\epsilon_T(w)+T^{-1}.
\end{aligned}
\]
Because the selected candidate minimizes this discrepancy,
\[
\max_{0\le k\le6}
|\widehat m_k(w)-\nu_*P_*^kr_*|
\le\epsilon_T(w)+T^{-1}.
\]
Applying the triangle inequality once more gives
\[
\begin{aligned}
\max_{0\le k\le6}|\nu_*P_*^kr_*-m_k|
&\le
\max_{0\le k\le6}|\nu_*P_*^kr_*-\widehat m_k(w)|
+\epsilon_T(w)\\
&\le2\epsilon_T(w)+T^{-1}.
\end{aligned}
\]
Combining this with the value modulus obtained above yields
\[
|\widehat\theta_T(w)-\theta|
\le
B_\alpha(2\epsilon_T(w)+T^{-1}).
\]
The right-hand side is nonnegative, so squaring preserves the inequality. Moreover,
\((2\epsilon_T(w)-T^{-1})^2\ge0\) gives
\[
\begin{aligned}
(\widehat\theta_T(w)-\theta)^2
&\le
B_\alpha^2(2\epsilon_T(w)+T^{-1})^2\\
&\le
B_\alpha^2\left(8\epsilon_T(w)^2+2T^{-2}\right).
\end{aligned}
\]
% lean: CausalSmith.Stat.PomdpBinaryhiddenRate.selected_grid_moment_error
% lean: CausalSmith.Stat.PomdpBinaryhiddenRate.selected_grid_intervention_moment_error
% lean: CausalSmith.Stat.PomdpBinaryhiddenRate.stableGridEstimator_value_error
% lean: CausalSmith.Stat.PomdpBinaryhiddenRate.stableGrid_upper

\item The sampling and grid terms admit separate horizon bounds. Since
\(V_{\alpha,L}\ge0\) and \(T\le2n_T\),
\[
56V_{\alpha,L}T
\le112V_{\alpha,L}n_T.
\]
Dividing by the positive product \(n_TT\) gives
\[
8\left(\frac{7V_{\alpha,L}}{n_T}\right)
=
\frac{56V_{\alpha,L}}{n_T}
\le
\frac{112V_{\alpha,L}}T.
\]
Also, \(T\ge1\) implies \(T^2\ge T>0\), and therefore
\[
T^{-2}\le T^{-1}.
\]
% lean: CausalSmith.Stat.PomdpBinaryhiddenRate.stableGrid_upper

\item The pointwise upper bound is integrable by integrability of \(\epsilon_T^2\). Monotonicity of integration for the nonnegative squared loss, followed by linearity of expectation under a probability law, now gives
\[
\begin{aligned}
\mathbb E_M[(\widehat\theta_T-\theta)^2]
&\le
\mathbb E_M\!\left[
 B_\alpha^2(8\epsilon_T^2+2T^{-2})
\right]\\
&=
B_\alpha^2\left(8\mathbb E_M[\epsilon_T^2]+2T^{-2}\right)\\
&\le
B_\alpha^2\left(
 8\frac{7V_{\alpha,L}}{n_T}+2T^{-2}
\right)\\
&\le
B_\alpha^2\left(
 \frac{112V_{\alpha,L}}T+\frac2T
\right)\\
&=
\frac{B_\alpha^2(112V_{\alpha,L}+2)}T.
\end{aligned}
\]
The fixed experiment was arbitrary, and every bound uses the same constants throughout the stated class.
% lean: CausalSmith.Stat.PomdpBinaryhiddenRate.stableGrid_upper
\end{enumerate}
\end{proof}

\begin{proof}[Proof of \cref{obj:thm:fixed-binary-minimax}]
\begin{enumerate}
\item Define the two testing experiments, their target values, and their observed laws by
\[
\begin{aligned}
\mathsf M_\pm
&:=\text{the experiment of \cref{obj:def:four-state-testing-family}
with }v=\pm v_T,\\
\theta_\pm&:=\theta(K_{\pm v_T}^{(T)},e),\\
\mathbb P_\pm^{\mathsf O_T}
&:=\mathcal L_{\mathsf M_\pm}(\mathsf O_T).
\end{aligned}
\]
By \cref{obj:lem:testing-family-membership},
\[
\mathsf M_\pm\in\mathcal M_T^{(2)}(t_0,\zeta),
\qquad
|\theta_+-\theta_-|=\frac1{8\sqrt T},
\qquad
\operatorname{KL}\!\left(
\mathbb P_+^{\mathsf O_T}\Vert\mathbb P_-^{\mathsf O_T}
\right)\leq\frac1{12},
\]
and the relative entropy is finite. Both experiments supply the same fair behavior and target policies to the estimator.
% lean: CausalSmith.Stat.PomdpBinaryhiddenRate.testingFamily_membership

\item Fix an admissible estimator \(\widetilde\theta\) from
\cref{obj:def:minimax-risk}. Because its supplied policies agree across the pair, it is the same measurable function of the observed trajectory in both experiments. Pinsker's inequality gives
\[
\left\|\mathbb P_+^{\mathsf O_T}
-\mathbb P_-^{\mathsf O_T}\right\|_{\mathrm{TV}}
\leq
\sqrt{\frac{
\operatorname{KL}(\mathbb P_+^{\mathsf O_T}
\Vert\mathbb P_-^{\mathsf O_T})}{2}}
\leq\sqrt{\frac1{24}}\leq\frac14.
\]
Define the error threshold and the two measurable error events by
\[
\varepsilon_{\mathrm{test}}:=\frac1{16\sqrt T},
\qquad
\mathcal E_\pm:=
\left\{|\widetilde\theta-\theta_\pm|
\geq\varepsilon_{\mathrm{test}}\right\}.
\]
Then
\[
|\theta_+-\theta_-|=2\varepsilon_{\mathrm{test}},
\qquad
\mathcal E_+^{\mathsf c}\subseteq\mathcal E_-,
\]
where the inclusion follows from the triangle inequality. Consequently,
\[
\begin{aligned}
\mathbb P_+^{\mathsf O_T}(\mathcal E_+)
+\mathbb P_-^{\mathsf O_T}(\mathcal E_-)
&\geq
\mathbb P_+^{\mathsf O_T}(\mathcal E_+)
+\mathbb P_-^{\mathsf O_T}(\mathcal E_+^{\mathsf c})\\
&\geq
1-\left\|\mathbb P_+^{\mathsf O_T}
-\mathbb P_-^{\mathsf O_T}\right\|_{\mathrm{TV}}
\geq\frac34.
\end{aligned}
\]
Thus
\[
\max\left\{
\mathbb P_+^{\mathsf O_T}(\mathcal E_+),
\mathbb P_-^{\mathsf O_T}(\mathcal E_-)
\right\}\geq\frac38.
\]
For each sign, integrating the pointwise inequality
\[
(\widetilde\theta-\theta_\pm)^2
\geq\varepsilon_{\mathrm{test}}^2\mathbf1_{\mathcal E_\pm}
\]
gives
\[
\mathbb E_{\mathsf M_\pm}
[(\widetilde\theta-\theta_\pm)^2]
\geq
\varepsilon_{\mathrm{test}}^2
\mathbb P_\pm^{\mathsf O_T}(\mathcal E_\pm).
\]
If the positive-sign error probability is at most the negative-sign error probability, apply this inequality to the negative sign; in the opposite case, apply it to the positive sign. In either case,
\[
\begin{aligned}
\max\left\{
\mathbb E_{\mathsf M_+}[(\widetilde\theta-\theta_+)^2],
\mathbb E_{\mathsf M_-}[(\widetilde\theta-\theta_-)^2]
\right\}
&\geq\frac38\varepsilon_{\mathrm{test}}^2\\
&=\frac3{2048T}
\geq\frac1{1024T},
\end{aligned}
\]
using \((\sqrt T)^2=T>0\).
% lean: CausalSmith.Stat.PomdpBinaryhiddenRate.testingPair_observedRisk_lower

\item For any admissible experiment \(\mathsf M\), define its risk by
\[
\mathcal R(\widetilde\theta,\mathsf M)
:=
\mathbb E_{\mathsf M}[(\widetilde\theta-\theta)^2],
\qquad
\mathsf M\in\mathcal M_T^{(2)}(t_0,\zeta).
\]
The target value in \cref{obj:def:target-value} is a stationary average of rewards in \([0,1]\). Since the estimator also takes values in \([0,1]\),
\[
0\leq\theta\leq1,
\qquad
0\leq(\widetilde\theta-\theta)^2\leq1,
\qquad
0\leq\mathcal R(\widetilde\theta,\mathsf M)\leq1.
\]
Thus every estimator's risks are bounded above over the class. The estimator family is nonempty, as witnessed by the constant estimator
\[
\widetilde\theta_0\equiv0.
\]
Membership of the testing pair and the preceding two-point bound therefore give
\[
\begin{aligned}
R_T^{(2)}(t_0,\zeta)
&=\inf_{\widetilde\theta}
\sup_{\mathsf M\in\mathcal M_T^{(2)}(t_0,\zeta)}
\mathcal R(\widetilde\theta,\mathsf M)\\
&\geq
\inf_{\widetilde\theta}
\max\left\{
\mathcal R(\widetilde\theta,\mathsf M_+),
\mathcal R(\widetilde\theta,\mathsf M_-)
\right\}
\geq\frac1{1024T}.
\end{aligned}
\]
% lean: CausalSmith.Stat.PomdpBinaryhiddenRate.observedRisk_unit_bounds
% lean: Causalean.Stat.le_minimaxValue_of_two_point

\item Use the stable-grid estimator of
\cref{obj:def:stable-grid-estimator}. Its selected reward vector and stationary probability vector satisfy
\[
0\leq r_*\leq\mathbf1,
\qquad
\pi(P_*)\geq0,
\qquad
\pi(P_*)\mathbf1=1.
\]
Hence
\[
0\leq\widehat\theta_T
=\pi(P_*)r_*
\leq\pi(P_*)\mathbf1=1.
\]
The finite candidate selection uses measurable comparison events, and the selected stationary reward is therefore measurable for each pair of supplied policies. This makes \(\widehat\theta_T\) admissible in
\cref{obj:def:minimax-risk}. By \cref{obj:thm:stable-grid-upper},
\[
\mathcal R(\widehat\theta_T,\mathsf M)
\leq\frac{B_\alpha^2(112V_{\alpha,L}+2)}{T}
\qquad
\text{for every }\mathsf M\in\mathcal M_T^{(2)}(t_0,\zeta).
\]
The class is nonempty because it contains \(\mathsf M_+\). Nonnegativity of the risks and the defining infimum and supremum now yield
\[
R_T^{(2)}(t_0,\zeta)
\leq
\sup_{\mathsf M\in\mathcal M_T^{(2)}(t_0,\zeta)}
\mathcal R(\widehat\theta_T,\mathsf M)
\leq\frac{B_\alpha^2(112V_{\alpha,L}+2)}{T}.
\]
% lean: CausalSmith.Stat.PomdpBinaryhiddenRate.stableGridEstimator_admissible
% lean: CausalSmith.Stat.PomdpBinaryhiddenRate.stableGrid_upper
% lean: Causalean.Stat.minimaxValue_le_worstCaseRisk_of_nonneg
% lean: Causalean.Stat.worstCaseRisk_le

\item The remaining assertions about the binary testing pair follow from
\cref{obj:lem:testing-family-membership}:
\[
e(a\mid x)=b(a\mid x)=\frac12,
\qquad d_e=d_b,
\]
and, for every \(C>1\),
\[
d_e(s)\leq C\,d_b(s)
\qquad\text{for every }s\in\mathcal S.
\]
Together with the estimator-wise two-point risk bound, these establish all the stated properties of the pair.
% lean: CausalSmith.Stat.PomdpBinaryhiddenRate.fixedBinary_minimax

\item Suppose now that \(\zeta>0\), and fix \(C>1\). Use the theorem's notation
\[
\beta=\frac2{2+t_0\zeta},
\qquad
q_C=\frac{C-1}{C}.
\]
The cited cardinality-uniform result supplies its signed-depth lower construction
\citep[Theorem 1 (Uniform overlap frontier) and the Signed-depth family definition]{CausalSmithLatentOverlap2026}.
Specifically, there are
\[
c_{\mathrm{Reset}}>0,\qquad
Q:\mathbb N\to\mathbb N,\qquad a_1,a_2>0
\]
such that, for every sufficiently large \(n\),
\[
a_1\log n\leq Q(n)\leq a_2\log n,
\qquad n\geq1,\qquad Q(n)\geq1.
\]
For these \(n\), define
\[
\begin{aligned}
\mathsf M_{n,\pm}^{\mathrm{Reset}}
&:=\text{the signed depth-}Q(n)\text{ reset experiments},\\
\theta_{n,\pm}^{\mathrm{Reset}}
&:=\text{their respective stationary target values}.
\end{aligned}
\]
Each experiment has exactly \(2(Q(n)+1)\) hidden states. For every admissible observable-data estimator
\(\widetilde\theta_{\mathrm{Reset}}\) in the reset decision problem, with values in \([-1,1]\), the cited construction gives
\[
\max\left\{
\mathbb E_{\mathsf M_{n,-}^{\mathrm{Reset}}}
[(\widetilde\theta_{\mathrm{Reset}}
-\theta_{n,-}^{\mathrm{Reset}})^2],
\mathbb E_{\mathsf M_{n,+}^{\mathrm{Reset}}}
[(\widetilde\theta_{\mathrm{Reset}}
-\theta_{n,+}^{\mathrm{Reset}})^2]
\right\}
\geq c_{\mathrm{Reset}}n^{-\beta}.
\]
Finally, the same eventual condition \(Q(n)\geq1\) gives
\[
2(Q(n)+1)\geq4>2,
\]
which proves the asserted cardinality inequality for the negative-sign experiment.
% lean: CausalSmith.Stat.PomdpBinaryhiddenRate.comparator_reset_depth_growth

\item Since \(t_0\zeta>0\),
\[
2+t_0\zeta>2>0,
\qquad
0<\beta=\frac2{2+t_0\zeta}<1,
\qquad
1-\beta>0.
\]
Also, \(C>1\) implies
\[
q_C>0,
\qquad q_C^{\,2(1-\beta)}>0.
\]
% lean: CausalSmith.Stat.PomdpBinaryhiddenRate.fixedBinary_minimax

\item Along the positive natural numbers, the negative-power limit gives
\[
n^{-(1-\beta)}\longrightarrow0.
\]
Multiplication by the fixed constant
\(q_C^{-2(1-\beta)}\) preserves this limit. For every \(n\geq1\), the laws of real powers give
\[
\frac{n^{-1}}{n^{-\beta}}
=n^{\beta-1}=n^{-(1-\beta)},
\]
and hence
\[
\frac{n^{-1}}
{n^{-\beta}q_C^{\,2(1-\beta)}}
=
n^{-(1-\beta)}q_C^{-2(1-\beta)}
\longrightarrow0.
\]
This proves the claimed limit and completes the proof.
% lean: CausalSmith.Stat.PomdpBinaryhiddenRate.fixedBinary_minimax
\end{enumerate}
\end{proof}

\begin{proof}[Proof of \cref{obj:prop:coincident-policy-reduction}]
We use epochs \(1,\ldots,T\), shifting the zero-based epochs of \cref{obj:lem:observable-intervention-moments} by one. Every sum over an empty range is understood to be zero.

\begin{enumerate}
\item Since \(e=b\), substitution in the policy-induced transition matrices gives, for every \(s=(x,h)\in\mathcal S\) and \(s'\in\mathcal S\),
\[
P_e(s,s')
=
\sum_{a\in\{0,1\}}e(a\mid x)
   \int_{[0,1]}K(dy,s'\mid s,a)
=
\sum_{a\in\{0,1\}}b(a\mid x)
   \int_{[0,1]}K(dy,s'\mid s,a)
=
P_b(s,s').
\]
The stationary-law assignment evaluated at these identical matrices likewise gives
\[
d_e=d_b.
\]
% lean: CausalSmith.Stat.PomdpBinaryhiddenRate.coincidentPolicy_kernel_stationary

\item We first verify membership in the class of \cref{obj:def:binary-pomdp-class} at logarithmic overlap parameter \(\zeta=0\). Sequential assignment and \(e=b\) give
\[
e(a\mid x)=b(a\mid x)\geq0,
\qquad
\sum_{a\in\{0,1\}}e(a\mid x)=1,
\qquad
e(a\mid x)\leq \exp(0)b(a\mid x)=b(a\mid x).
\]
Thus \cref{obj:ass:policy-overlap} holds with \(L=1\). Moreover, for every \(\mu,\mu'\in\Delta(\mathcal S)\),
\[
\|\mu P_e-\mu'P_e\|_{\mathrm{TV}}
=
\|\mu P_b-\mu'P_b\|_{\mathrm{TV}}
\leq
\alpha\|\mu-\mu'\|_{\mathrm{TV}},
\]
so \cref{obj:ass:target-contraction} follows from \cref{obj:ass:behavior-contraction}. Together with the remaining hypotheses, these facts establish membership in \(\mathcal M_T^{(2)}(t_0,0)\).

Write the target-policy conditional mean reward vector as
\[
r_e(s):=
\sum_{a\in\{0,1\}}e(a\mid x)
\sum_{s'\in\mathcal S}\int_{[0,1]}y\,K(dy,s'\mid s,a),
\qquad s=(x,h)\in\mathcal S.
\]
Applying \cref{obj:lem:observable-intervention-moments} at depth zero gives, at every observed epoch,
\[
m_0=d_bP_e^0r_e=d_br_e=d_er_e=\theta,
\]
where the last equality is \cref{obj:def:target-value}.
% lean: CausalSmith.Stat.PomdpBinaryhiddenRate.observable_intervention_moment_zero

\item To identify this moment with the ordinary reward mean, define the observed action likelihood ratio and depth-zero score by
\[
\rho_t:=
\begin{cases}
e(A_t\mid X_t)/b(A_t\mid X_t),& b(A_t\mid X_t)>0,\\
0,& b(A_t\mid X_t)=0,
\end{cases}
\qquad
Z_t(0):=Y_t\rho_t,
\qquad 1\leq t\leq T.
\]
Here probability and expectation are taken under the behavior experiment. By \cref{obj:ass:sequential-ignorability},
\[
\Pr_b\!\left(b(A_t\mid X_t)=0\right)
=
\mathbb E_b\!\left[
\sum_{\substack{a\in\{0,1\}\\ b(a\mid X_t)=0}}
b(a\mid X_t)
\right]
=0.
\]
On the complementary event, \(e=b\) gives \(\rho_t=1\). Consequently,
\[
Z_t(0)=Y_t\quad\text{almost surely},
\qquad
\mathbb E_b[Y_t]
=
\mathbb E_b[Z_t(0)]
=
m_0
=
\theta.
\]
% lean: CausalSmith.Stat.PomdpBinaryhiddenRate.coincidentPolicy_reward_mean

\item Define the sample mean by
\[
\overline Y_T:=T^{-1}\sum_{t=1}^{T}Y_t.
\]
The bounded rewards are integrable, so linearity of expectation and \(T\geq1\) yield
\[
\mathbb E_b[\overline Y_T]
=
T^{-1}\sum_{t=1}^{T}\mathbb E_b[Y_t]
=
T^{-1}T\theta
=
\theta.
\]
% lean: CausalSmith.Stat.PomdpBinaryhiddenRate.coincidentPolicy_sampleMean_unbiased

\item The reward support and \cref{obj:ass:pomdp-kernel} give
\[
\Pr_b(Y_t\in[0,1])
=
\mathbb E_b\!\left[
K([0,1]\times\mathcal S\mid S_t,A_t)
\right]
=1.
\]
Thus every reward is square integrable. For square-integrable real random variables \(W_1,W_2\) under the behavior law, use
\[
\operatorname{Cov}_b(W_1,W_2)
:=
\mathbb E_b\!\left[
(W_1-\mathbb E_b[W_1])(W_2-\mathbb E_b[W_2])
\right],
\qquad
\operatorname{Var}_b(W_1)
:=
\operatorname{Cov}_b(W_1,W_1).
\]
Since \(t_0>0\),
\[
0\leq\alpha=\exp(-1/t_0)<1.
\]
Also \(Y_t^2\leq1\) almost surely, and hence
\[
\mathbb E_b[Y_t^2]\leq\mathbb E_b[1]=1,
\qquad
\operatorname{Cov}_b(Y_t,Y_t)
=
\mathbb E_b[Y_t^2]-(\mathbb E_b[Y_t])^2
\leq1.
\]
% lean: CausalSmith.Stat.PomdpBinaryhiddenRate.coincidentPolicy_sampleMean_mse

\item Expanding the variance of the reward sum gives
\[
\operatorname{Var}_b\!\left(\sum_{t=1}^{T}Y_t\right)
=
\sum_{t=1}^{T}\sum_{u=1}^{T}
\operatorname{Cov}_b(Y_t,Y_u).
\]
For \(1\leq t<u\leq T\), \cref{obj:lem:observable-intervention-moments}, applied at depth zero in the class established above, supplies
\[
\left|
\mathbb E_b\!\left[
(Z_t(0)-m_0)(Z_u(0)-m_0)
\right]
\right|
\leq\alpha^{u-t-1}.
\]
Using the almost-sure score identities and reward means already established,
\[
\operatorname{Cov}_b(Y_t,Y_u)
=
\mathbb E_b\!\left[
(Z_t(0)-m_0)(Z_u(0)-m_0)
\right],
\]
and therefore
\[
\operatorname{Cov}_b(Y_t,Y_u)
\leq
\left|\operatorname{Cov}_b(Y_t,Y_u)\right|
\leq\alpha^{u-t-1}.
\]
% lean: CausalSmith.Stat.PomdpBinaryhiddenRate.coincidentPolicy_reward_covariance

To bound each covariance row, first use the finite geometric identity
\[
(1-\alpha)\sum_{\ell=0}^{T-1}\alpha^\ell
=
1-\alpha^T
\leq1,
\qquad
\sum_{\ell=0}^{T-1}\alpha^\ell
\leq\frac1{1-\alpha}.
\]
Here \(\ell\) is a nonnegative integer exponent. For fixed \(t\), the exponent maps \(u\mapsto t-u-1\) on \(1\leq u<t\) and \(u\mapsto u-t-1\) on \(t<u\leq T\) are injective and take values in \(\{0,\ldots,T-1\}\). Because all powers of \(\alpha\) are nonnegative, symmetry of covariance gives
\[
\begin{aligned}
\sum_{u=1}^{t-1}\operatorname{Cov}_b(Y_t,Y_u)
&=
\sum_{u=1}^{t-1}\operatorname{Cov}_b(Y_u,Y_t)
\leq
\sum_{u=1}^{t-1}\alpha^{t-u-1}
\leq
\sum_{\ell=0}^{T-1}\alpha^\ell
\leq\frac1{1-\alpha},\\
\sum_{u=t+1}^{T}\operatorname{Cov}_b(Y_t,Y_u)
&\leq
\sum_{u=t+1}^{T}\alpha^{u-t-1}
\leq
\sum_{\ell=0}^{T-1}\alpha^\ell
\leq\frac1{1-\alpha}.
\end{aligned}
\]
These inequalities include the empty past or future ranges at the endpoints. Splitting the row into its diagonal, past, and future parts now yields
\[
\begin{aligned}
\sum_{u=1}^{T}\operatorname{Cov}_b(Y_t,Y_u)
&=
\operatorname{Cov}_b(Y_t,Y_t)
+\sum_{u=1}^{t-1}\operatorname{Cov}_b(Y_t,Y_u)
+\sum_{u=t+1}^{T}\operatorname{Cov}_b(Y_t,Y_u)\\
&\leq1+\frac2{1-\alpha}.
\end{aligned}
\]
Summing over \(t\) proves
\[
\operatorname{Var}_b\!\left(\sum_{t=1}^{T}Y_t\right)
\leq
T\left(1+\frac2{1-\alpha}\right).
\]
% lean: CausalSmith.Stat.PomdpBinaryhiddenRate.geometricCovariance_row_sum_le

\item Finally, unbiasedness identifies the mean squared error with the variance. Rescaling the preceding bound, using \(T>0\), gives
\[
\begin{aligned}
\mathbb E_b\!\left[(\overline Y_T-\theta)^2\right]
&=
\operatorname{Var}_b(\overline Y_T)\\
&=
T^{-2}\operatorname{Var}_b\!\left(\sum_{t=1}^{T}Y_t\right)\\
&\leq
T^{-2}T\left(1+\frac2{1-\alpha}\right)
=
\frac{1+2/(1-\alpha)}{T}.
\end{aligned}
\]
% lean: CausalSmith.Stat.PomdpBinaryhiddenRate.coincidentPolicy_sampleMean_mse
\end{enumerate}
\end{proof}

\begin{proof}[Proof of \cref{obj:prop:exhaustive-grid-arithmetic-complexity}]
We count integer-coordinate codes throughout, including at resolution zero.
\begin{enumerate}
\item For \(M\in\mathbb Z_{\ge0}\), define the integer simplex grid by
\[
\mathcal I_M
:=
\left\{
(u_1,u_2,u_3,u_4)\in\{0,\ldots,M\}^{4}:
\sum_{i=1}^{4}u_i=M
\right\}.
\]
Taking coordinate values gives a bijection between \(\mathcal I_M\) and the nonnegative integer quadruples summing to \(M\): each coordinate of such a quadruple is at most its sum, hence at most \(M\). These quadruples are in turn in bijection with multisets of size \(M\) on the four labels \(1,\ldots,4\), by assigning multiplicity \(u_i\) to label \(i\). The multiset counting formula and binomial symmetry therefore give
\[
|\mathcal I_M|
=
\binom{M+3}{M}
=
\binom{M+3}{3}.
\]
This also covers \(M=0\), where the unique quadruple is \((0,0,0,0)\).
% lean: CausalSmith.Stat.PomdpBinaryhiddenRate.simplexGrid_card

\item Define the transition-row codes, reward codes, and unfiltered candidate codes by
\[
\mathcal J_M:=\mathcal I_M^{4},
\qquad
\mathcal E_M:=\{0,\ldots,M\}^{4},
\qquad
\mathcal U_M:=\mathcal I_M\times
   (\mathcal J_M\times\mathcal E_M).
\]
A choice of four rows determines a unique transition array, and its rows recover the choice. Thus this correspondence is injective, and the product counting formula yields
\[
|\mathcal J_M|
=
\prod_{i=1}^{4}|\mathcal I_M|
=
|\mathcal I_M|^{4}.
\]
% lean: CausalSmith.Stat.PomdpBinaryhiddenRate.gridMatrixRows_card
Each reward coordinate has \(M+1\) choices, so
\[
|\mathcal E_M|=(M+1)^4.
\]
The candidate enumeration is precisely the displayed Cartesian product: a triple determines, and is determined by, its initial-vector code, transition-row code, and reward code. Consequently,
\[
|\mathcal U_M|
=
|\mathcal I_M|\,|\mathcal J_M|\,|\mathcal E_M|
=
|\mathcal I_M|^{5}(M+1)^4
=
\binom{M+3}{3}^{5}(M+1)^4.
\]
For \(M\ge1\), coordinatewise division by \(M\) gives the unfiltered grid triples.
% lean: CausalSmith.Stat.PomdpBinaryhiddenRate.gridCandidateUniverse_card

\item For a transition-row code
\(\mathsf R=(\mathsf R_{ij})\in\mathcal J_M\), define its decoded transition candidate by
\[
R_{ij}:=\frac{\mathsf R_{ij}}{M},
\qquad 1\le i,j\le4,
\]
with division by zero interpreted as zero. The contraction coefficient used for filtering is
\[
\delta(R)
:=
\max_{1\le i,j\le4}
\frac12\sum_{\ell=1}^{4}|R_{i\ell}-R_{j\ell}|.
\]
For every \(M\in\mathbb Z_{\ge0}\) and every \(\alpha\in\mathbb R\), define the retained codes by
\[
\mathcal U_M^\alpha
:=
\left\{
(u,\mathsf R,w)\in\mathcal U_M:
\delta(R)\le\alpha+\frac6M
\right\},
\]
using the same division convention. For \(M\ge1\), coordinatewise division by \(M\) identifies the retained codes \(\mathcal U_M^\alpha\) with the candidate set in \cref{obj:def:stable-grid-estimator}. This is a subset of \(\mathcal U_M\), so
\[
|\mathcal U_M^\alpha|
\le
|\mathcal U_M|
=
\binom{M+3}{3}^{5}(M+1)^4.
\]
In particular, at \(M=0\) the decoded transition array is zero and the subset argument still applies.
% lean: Finset.card_filter_le

\item The standard binomial bound
\[
\binom{M+3}{3}\le(M+3)^3
\]
together with \(M+1\le M+3\) gives
\[
\begin{aligned}
|\mathcal U_M|
&=\binom{M+3}{3}^{5}(M+1)^4\\
&\le \bigl((M+3)^3\bigr)^5(M+3)^4\\
&=(M+3)^{19}.
\end{aligned}
\]
% lean: CausalSmith.Stat.PomdpBinaryhiddenRate.gridCandidateUniverse_growth

\item Now fix the mixing-scale parameter \(t_0>0\) and its associated contraction bound. Then
\[
\alpha=\exp(-1/t_0)>0.
\]
Define
\[
c:=(26+36/\alpha)^{19}>0.
\]
For an integer trajectory length \(T\ge1\), take
\[
M:=\left\lceil(22+36/\alpha)T\right\rceil.
\]
Since \((22+36/\alpha)T\ge0\), the ceiling bound gives
\[
M<(22+36/\alpha)T+1.
\]
Adding three and using \(T\ge1\), we obtain
\[
M+3
<
(22+36/\alpha)T+4
\le
(26+36/\alpha)T,
\]
and hence
\[
M+3\le(26+36/\alpha)T.
\]
% lean: CausalSmith.Stat.PomdpBinaryhiddenRate.gridSize_linear_bound
Both sides are nonnegative, so raising this inequality to the nineteenth power and applying the preceding cardinality bound yields
\[
|\mathcal U_M|
\le
(M+3)^{19}
\le
\bigl((26+36/\alpha)T\bigr)^{19}
=
cT^{19}.
\]
Finally, the filtering inequality gives
\[
|\mathcal U_M^\alpha|
\le
|\mathcal U_M|
\le
cT^{19}.
\]
The chosen constant depends on \(t_0\) through \(\alpha\) and is valid for every integer \(T\ge1\).
% lean: CausalSmith.Stat.PomdpBinaryhiddenRate.exhaustiveGrid_candidate_count
\end{enumerate}
\end{proof}
